\section{Results}
\label{sec:results}
\label{sec:experiments}

\subsection{Why the gate was stopped}

The registered Gate~3 pilot read \emph{kill-or-reframe} on two triggers. The first was
headroom: even a clairvoyant that knows the shock at onset clears the registered 5\% profit bar
on only two of six families (a demand rise and a lead-time shift), because pulses, loss bursts
and the compound family's transit pauses are over before an order placed at onset can land, and
a demand fall offers little to recover. The second compared the gate with a detector compiled
without language and found it behind by 566 per episode in gross profit; forensic replays traced
that gap to the endpoint, since on net reward the same contrast is $+2{,}306$. The exploratory
pilot then gave the gate real model content. Gemini and Grok named the right family and
direction in 59 and 53 of 96 shocked episodes, yet the gate
activated only 43\% and 39\% of the right-family proposals, a median of 10 and 11 periods after
the proposal, and ended at $-234$ and $-172$ net per episode against arm~1 with intervals
including zero. The magnitude field was right in 8 and 7 of 96 episodes, partly because the
alert templates themselves label every non-distractor shock \emph{medium}. Acting at once was
far worse: arm~8 lost 2{,}549 and 2{,}624 per episode.

\subsection{The fresh-seed confirmation}

Table~\ref{tab:confirm} gives the primary contrasts on the 240 fresh episodes. Certify-then-hedge
beats the gate by $309$ (Gemini) and $179$ (Grok) net per episode and arm~1 by $178$ and $163$;
all four one-sided tests reject after Holm correction (largest $p = 0.0031$). In a post hoc check
without the six clusters that a structural test had touched before registration, the four
contrasts are $+329$, $+188$, $+163$ and $+145$ (largest $p = 0.0024$). Gross profit agrees on
both primary contrasts. On the 48 null episodes the method is non-inferior to arm~1 at the
registered margin of $-25$ with both models (Gemini $-3.2$ $[-8.2, -0.1]$, Grok $+1.9$
$[-1.7, 8.2]$). With Gemini the small loss is distinguishable from zero; post hoc, it comes from
the 24 false-alert twins ($-6.4$ $[-16.1, -0.3]$, worst episode $-105$), since silent twins
rarely trigger. The gate itself is no better than arm~1 on fresh seeds ($-131$ and $-16$ net,
intervals including zero), and acting at once loses $2{,}046$ and $2{,}284$.
Fig.~\ref{fig:models} repeats the confirmation with four further proposer models: every interval
for the method against the gate and against arm~1 excludes zero.

\begin{table}[!tb]
\caption{Fresh-seed confirmation (registered; exploratory with respect to the project's
preregistration), 240 episodes, 48 seed clusters, $p/h = 4$. Mean paired difference per
episode, net reward, with the 95\% seed-cluster bootstrap interval; gross profit in the last
column. $^\ast$Primary tests, all rejected after Holm correction.}
\label{tab:confirm}
\centering
\footnotesize
\setlength{\tabcolsep}{3pt}
\begin{tabular}{@{}llrr@{}}
\toprule
Contrast & Model & Net $[95\%]$ & Gross \\
\midrule
hedge $-$ gate$^\ast$ & Gemini & $+309\ [173, 472]$ & $+85$ \\
                      & Grok   & $+179\ [87, 281]$  & $+139$ \\
hedge $-$ arm 1$^\ast$ & Gemini & $+178\ [55, 326]$ & $+280$ \\
                       & Grok   & $+163\ [48, 300]$ & $+262$ \\
gate $-$ arm 1         & Gemini & $-131\ [-316, 28]$ & $+195$ \\
                       & Grok   & $-16\ [-104, 75]$  & $+123$ \\
at once $-$ arm 1      & Gemini & $-2{,}046\ [-2{,}612, -1{,}478]$ & $+501$ \\
                       & Grok   & $-2{,}284\ [-2{,}849, -1{,}710]$ & $+501$ \\
content-free $-$ arm 1 & either & $+169\ [48, 312]$ & $+270$ \\
hedge $-$ content-free & Gemini & $+9\ [-5, 24]$ & $+10$ \\
                       & Grok   & $-6\ [-29, 11]$ & $-8$ \\
\bottomrule
\end{tabular}
\end{table}

The gross column shows the endpoint trap of our first registration in its plainest form. Among
the arms in Table~\ref{tab:confirm}, acting at once looks \emph{best} on gross profit ($+501$),
because gross profit charges no holding and so pays for over-stocking, while on net it is the
worst by a wide margin; the hedge's contrast with acting at once changes sign too. Every claim
here is on net.

\begin{figure*}[!t]
\centering
\includegraphics[width=\textwidth]{fig2_models.pdf}
\caption{The fresh-seed confirmation with six proposer models (240 episodes each).
\textbf{(a)}~First proposals naming the right family and direction, among the 170 shocked
episodes where the trigger fired. \textbf{(b)}~Acting on proposals at once (arm~8) against
arm~1, net per episode. \textbf{(c)}~The gate (arm~10) and certify-then-hedge against arm~1,
with 95\% seed-cluster intervals; the band and dashed line are the content-free control
($\lambda = 1$, no model call, the same for every model). Perception and the harm of acting at
once vary widely with the model; the hedge's gain does not, and it stays within the band.}
\label{fig:models}
\end{figure*}

\subsection{Where the gain comes from}

Against arm~1 the gain is where the headroom is, and it is concentrated: by the registered
per-family readout, the 32 shocked lead-time episodes supply 96\% of it. The method earns
$+1{,}286$ (Gemini) and $+1{,}177$ (Grok) per lead-time episode, against $+90$ and $+88$ per
demand fall and $+95$ and $+80$ per pulse; it gives back $83$ per demand rise, about $40$ per
compound and $3$ per loss episode. The following readouts are post hoc. Without the lead-time
units (their episodes and twins) the method is $+9$ $[-24, 43]$ and $+7$ $[-25, 36]$ against
arm~1. Against the gate the advantage is broader: without those units the method still leads by
$+227$ $[82, 408]$ with Gemini ($+58$ $[-1, 134]$ with Grok), because the gate compiles the
stated magnitude and loses $1{,}057$ per demand-rise episode with Gemini. Two mechanisms carry
the lead-time gain. The first is sizing: when a model named a lead-time shift it usually
oversized it (medium or high in 11 of 17 Gemini proposals and all 11 of Grok's, while every true
shift was one period), and the gate compiles medium as two periods; the within-hypothesis
posterior recovered the true offset from receipts (Fig.~\ref{fig:episode}). The second is
timing: the hedge commits partially as soon as evidence accumulates, where the gate commits all
at once a median of 11 periods after the proposal. Both are statistical, which the next result
makes explicit.

\subsection{What the language contributes}

The content-free variant calls no model but is timed by the same alerts. It earns $+169$
$[48, 312]$ against arm~1, and the method's net reward is not distinguishable from it ($+9$
$[-5, 24]$ with Gemini, $-6$ $[-29, 11]$ with Grok); with the four further models the difference
runs from $-43$ to $0$, with intervals below zero for Llama~3.1~8B and GPT-3.5 Turbo. At this
cost ratio, then, a certified hypothesis search timed by alerts recovers the headroom without
detectable help from the text. What the model's family label changes is commitment on null
episodes: the summed shock mass falls from $0.83$ (content-free) to $0.52$ with Gemini and
$0.46$ with Grok. Post hoc, these cuts of 37\% and 45\% (intervals $[33, 57]\%$ and
$[42, 61]\%$) fall almost entirely on false-alert twins, with no detectable change in null
reward ($-2$ $[-8, 2]$ and $+3$ $[-0.04, 8.4]$). Naming any family has this effect, right or
wrong, because it moves prior weight off the hypotheses the model does not name. Trusting the
family alone ($\lambda = 0$) cuts null commitment further, to $0.18$ and $0.06$, but costs net
reward against the content-free arm with Grok ($-83$ $[-169, -13]$) and Llama~3.1~8B ($-192$),
and is not distinguishable from it with Gemini ($-16$ $[-66, 27]$); hence the hedged prior.

\subsection{Cost ratio}

The confirmation fixed $p/h = 4$, the middle of InventoryBench's three ratios. Post hoc, we
replayed its 240 episodes at $p/h = 1$, $2$, $9$ and $19$, holding each model's answers at what it
said under $p/h = 4$. The prompt shows the margin, the inventory position and the order and
arrival history, all of which change with the ratio, so the model-calling arms are a
counterfactual with the answer fixed; the others are exact, and at $p/h = 4$ the replay
reproduces every registered record. The method beats the gate at every ratio with both models
($+85$ to $+722$, every interval above zero). Its gain over arm~1 is small when one period of holding costs the
whole margin ($p/h = 1$: $+16$ $[-14, 47]$ and $+10$ $[-18, 40]$) and grows with the margin, to
$+1{,}056$ $[398, 1{,}850]$ and $+991$ $[376, 1{,}739]$ at $p/h = 19$, where it no longer rests
on lead-time shifts alone ($+125$ $[43, 227]$ and $+121$ $[38, 224]$ without those units). Part
of that growth is scale, since every reward grows with $p$: as a share of arm~1's net reward the
gain is 0.5\% and 0.3\% at $p/h = 1$, 1.1\% and 1.0\% at 4, and 1.2\% at 19 with both models.
The text's share moves with the ratio too: with Gemini the method beats the content-free control
at $p/h = 9$ ($+47$ $[6, 97]$) and $19$ ($+84$ $[18, 166]$), while with Grok the label costs a
little at $p/h = 1$ ($-5$ $[-10, -0.02]$) and is not distinguishable from zero at the other
ratios.

\subsection{Model scale}

The ladder ran 20 of the 21 models on the 120 pilot episodes (StepFun's slowest reasoning model
had not finished at the time of writing). First proposals named the right family and direction
in 0 to 59 of 96 shocked episodes. Llama~3.2 1B never produced a non-abstaining \shockspec{},
so every proposal fell back to the baseline; 3B produced seven, all demand rises, on which
acting at once lost 28 per episode $[-54, -6]$; and Qwen~2.5 7B abstained or failed on almost
every call. Acting on proposals at once lost between nothing (Llama~3.2 1B) and $3{,}656$ net
per episode (Gemini~2.5 Flash-Lite, which almost never abstains). The loss grows with how
readily a model commits, and being right does not protect it: across the 20 models commitment
and accuracy move together (post hoc, $r = 0.86$), and the most accurate models lose among the
most. The gate stayed between $-300$ and $+3$ against arm~1 and acted on a false alert in at most
one of 12 null episodes; contrary to the stage's registered expectation, its interval lay below
zero on two models (Gemini~2.5 Flash $-294$ $[-567, -45]$, GPT-3.5 Turbo $-60$ $[-151, -1]$).
Certify-then-hedge's development replays (no claim) had point estimates between $+107$ and
$+158$ on every model, including those that produced nothing usable, against $+136$ for its
content-free variant.
